\documentclass[11pt]{article}
\usepackage[margin=1in]{geometry}
\usepackage{amsmath,amssymb,amsthm}
\newtheorem{theorem}{Theorem}
\newtheorem{corollary}{Corollary}
\title{A Rational Arm Certificate Covering Every Physical Angle}
\author{}\date{October 5, 2026 --- paper proof}
\begin{document}\maketitle
Use $m(z)=3z/2-1$ on $[0,1]$, $m(z)=z/2$ on $[1,2]$, and $m(z)=1$ on $[2,\infty)$.

\begin{theorem}[Uniform bootstrap up to $8/5$]
Let $0<L\le8/5$. Suppose continuous nonnegative functions $f,g$ on $[0,L]$ satisfy
\[
 f(t)\ge1+\int_0^t m(g(s))\,ds,\qquad
 g(t)\ge1+\int_t^L m(f(s))\,ds.
\]
Then $f(t)\ge1+t/2$ and $g(t)\ge1+(L-t)/2$.
\end{theorem}
\begin{proof}
We give an exact finite rational certificate, not a floating-point iteration. Subdivide the parameter interval into twelve equal cells. An integer row $B=(B_0,\ldots,B_{11})$, with $0\le B_i\le1000$, is a lower certificate if
\[
 f(t)\ge B_i/1000\quad(iL/12\le t\le(i+1)L/12),
\]
and $g$ satisfies the reflected bounds. The zero row is valid.

Given a valid row, put $H_0=15000$ and, for $1\le i\le12$,
\[
 H_i=15000-2000i+3\sum_{j=0}^{i-1}B_{11-j}.
\]
At length $R=8/5$, the lower bound obtained by integrating $m$ of the reflected step function has value $H_i/15000$ at the $i$th grid point. It is linear on each cell. At a shorter length $L$, its value at the same relative position is $1+(L/R)(H/15000-1)$, which is at least $\min\{1,H/15000\}$. Monotonicity of $m$ justifies replacing the actual functions by these lower step functions. Therefore the new row
\[
 B_i^{\rm new}=\min\left\{1000,\max\left\{B_i,0,
 \left\lfloor\frac{\min(H_i,H_{i+1})}{15}\right\rfloor\right\}\right\}
\]
is valid simultaneously for every $0<L\le8/5$. At shared cell endpoints both bounds follow by continuity.

Eight applications of this integer recurrence give the following table. Every entry is rounded down, never to nearest, and the recurrence above permits direct verification of every entry.
\[
\setlength{\arraycolsep}{3pt}\scriptsize
\begin{array}{c|rrrrrrrrrrrr}
 j\backslash i&0&1&2&3&4&5&6&7&8&9&10&11\\\hline
0&0&0&0&0&0&0&0&0&0&0&0&0\\
1&866&733&600&466&333&200&66&0&0&0&0&0\\
2&866&733&600&466&333&213&119&53&13&0&0&12\\
3&869&735&602&471&348&239&148&81&41&28&28&41\\
4&874&747&619&494&377&273&187&124&85&72&72&85\\
5&883&764&645&529&420&325&246&188&153&144&144&160\\
6&898&794&689&586&491&407&338&289&261&257&257&276\\
7&921&839&758&676&601&535&483&448&432&432&436&462\\
8&959&912&866&819&775&738&712&699&699&701&719&753
\end{array}
\]
Thus $f,g\ge699/1000>2/3$. Since $m(2/3)=0$, the integral inequalities give $f,g\ge1$. One further substitution, using $m(1)=1/2$, gives the asserted linear bounds.
\end{proof}

\begin{corollary}[Every fixed-angle maximizer admits the majorant]
For every $0<\omega<\pi/2$, every positive-area global cap maximizer satisfies
\[
 \mathcal A_\omega(K)\le\mathcal A_1(K).
\]
Its canonical corner is continuously differentiable and has strictly decreasing first coordinate.
\end{corollary}
\begin{proof}
The fixed-angle curvature theorem supplies the integral inequalities and continuous nonnegative arms, with $L=\omega$. Since $\pi/2<8/5$, the certificate applies. The corner derivative is $-(f-1)u_t+(g-1)v_t$, whose first coordinate is strictly negative in the interior. Apply the horizontal-monotonicity majorization theorem, which needs no prior fan-containment hypothesis.
\end{proof}

\paragraph{Scope and verification.}
This extends the elementary three-step proof for $L\le1$ in \texttt{arm-bootstrap.tex}; it is not needed for the already proved theorem on $[0,1]$. It supplies the structural majorization for larger fixed angles, not an explicit optimizer there. The table was checked using both integer arithmetic and Python's exact \texttt{Fraction} arithmetic. The finite recurrence and all certificate entries are part of the paper proof; the script is an optional reproducibility aid, not an unstated numerical hypothesis. No CI or Lean compilation was used.
\end{document}
